Automated symptom-based disease prediction systems hold significant potential for modern Clinical Decision Support Systems (CDSS). However, prevailing machine learning implementations on public benchmark datasets frequently report near-perfect ($>99\%$) or 100\% classification accuracies without examining the underlying data fidelity. In this research, an independent empirical forensic audit of the widely benchmarked Kaggle Disease Prediction dataset is conducted. The audit uncovers critical structural anomalies: although the training set comprises 4,920 records across 41 disease classes (disproving the published 42-class description), it contains only 304 unique symptom vectors, representing an extreme synthetic duplication rate of 93.8\%. Furthermore, 41 of the 42 official test records (97.6\%) are exact duplicate copies of training records, demonstrating severe train-test leakage that renders standard accuracy metrics clinically meaningless. To address these vulnerabilities, an uncertainty-aware evaluation framework is established across three distinct protocols: standard split (Protocol A), stratified cross-validation (Protocol B), and duplicate-aware group evaluation (Protocol C). Leakage-aware multi-method feature selection demonstrates that a compact subset of 50 consensus symptoms preserves $\ge 95\%$ operational Macro-F1 retention. Furthermore, controlled symptom sparsity simulations ($s \in [1, 10]$ presenting symptoms relative to the cohort natural mean of $7.45 \pm 3.59$) demonstrate that differential diagnostic candidate ranking (Top-3 and Top-5) provides a resilient decision-support safeguard when single-label Top-1 accuracy declines under sparse input vectors on this benchmark. Crucially, statistical stability across feature selection algorithms establishes dataset-level algorithmic evidence rather than independent clinical validation. Finally, an experimental roadmap is outlined for subsequent research phases, formulating planned methodologies for stochastic reporting noise modeling, empirical probability calibration, and TreeSHAP clinical literature alignment audits.

\vspace{8pt}
\textbf{Keywords:} Disease Prediction, Clinical Decision Support Systems (CDSS), Data Leakage, Duplicate-Aware Evaluation, Feature Selection, Symptom Sparsity, Differential Diagnosis.

